% !TeX program = xelatex
% UTF-8 transcription; source wording and highlighted points are preserved.
% images/pic01.jpg = 20261013_064508163.jpg (page 1)
% images/pic02.jpg = 20261013_064511934.jpg (page 2)
% Compile directly with XeLaTeX, or input into a XePersian document.
\newif\ifNephroticStandalone
\ifdefined\settextfont
  \NephroticStandalonefalse
\else
  \NephroticStandalonetrue
\fi
\ifNephroticStandalone
  \documentclass[11pt,a4paper]{article}
  \usepackage[margin=1.8cm]{geometry}
  \usepackage{amsmath,amssymb,xcolor,enumitem}
  \usepackage{xepersian}
  \IfFontExistsTF{Amiri}{\settextfont{Amiri}}{\settextfont{FreeSerif}}
  \setlatintextfont{FreeSerif}
  \begin{document}
\fi
\begingroup
\emergencystretch=2em
\ifdefined\setlist
  \setlist[itemize]{nosep,topsep=3pt}
\fi
\newcommand{\NSen}[1]{\lr{#1}}
\ifdefined\colorbox
  \newcommand{\NShl}[1]{\colorbox{yellow!35}{\strut #1}}
\else
  \newcommand{\NShl}[1]{\textbf{#1}}
\fi

\section*{سندرم نفروتیک ـ صفحهٔ ۱}
\subsection*{پروتئین‌اوری}
\begin{itemize}
\item پروتئین‌اوری نفروتیک در کودکان: دفع پروتئین
  \NSen{$>40\,\mathrm{mg}/\mathrm{m}^{2}/\mathrm{hour}$}
  یا \NSen{$U_{\mathrm{Pr}/\mathrm{Cr}}>2\,\mathrm{mg}/\mathrm{mg}$}.
\item الگوهای پروتئین‌اوری: \NSen{Transient}، ارتواستاتیک، \NSen{Persistent}.
\item سندرم نفروتیک یک پروتئین‌اوری \NSen{persistent} است.
\end{itemize}

\subsection*{\NSen{Transient}}
\begin{itemize}
\item در کودکان با \NShl{تب}، \NShl{دهیدریشن}، به دنبال \NShl{تشنج}، \NShl{ورزش شدید}.
\item ظرف چند روز رفع می‌شود.
\item \NShl{نیاز به ارزیابی اضافه ندارد}.
\end{itemize}

\subsection*{پروتئین‌اوری ارتواستاتیک}
\begin{itemize}
\item پروتئین‌اوری خفیف در زمان سرپا ایستادن.
\item \NSen{$U_{\mathrm{Pr}/\mathrm{Cr}}=0.2\text{--}1\,\mathrm{mg}/\mathrm{mg}$}.
\item در نوجوانان قدبلند و لاغر.
\item با بیماری پیشروندهٔ کلیه همراهی ندارد.
\end{itemize}

\subsection*{\NSen{Persistent}}
\begin{itemize}
\item منشأ گلومرولار یا توبولار.
\item گلومرولار:
  \begin{itemize}
  \item پروتئین با وزن‌های مولکولی \NShl{متفاوت}.
  \item معمولاً \NShl{هماچوری}، \NShl{\NSen{RBC} کست}، \NShl{\NSen{HTN}} وجود دارد.
  \end{itemize}
\item توبولار:
  \begin{itemize}
  \item پروتئین با وزن مولکولی \NShl{پایین} دفع می‌شود.
  \item ممکن است \NSen{dipstick} آن‌ها منفی باشد.
  \end{itemize}
\end{itemize}

\subsection*{\NShl{سندرم فانکونی}}
\begin{itemize}
\item \NShl{\NSen{RTA} تیپ ۲}، هایپوکالمی، فسفات‌اوری، آمینواسیداوری.
\item پروتئین‌اوری توبولار + توبولار الکترولیت وستینگ + \NShl{گلوکوزاوری}.
\item \NShl{پروگزیمال رنال توبولار اسیدوز}.
\end{itemize}

\subsection*{سندرم نفروتیک}
\begin{itemize}
\item پروتئین‌اوری بالای \NSen{$40\,\mathrm{mg}/\mathrm{m}^{2}/\mathrm{hour}$}.
\item هایپوآلبومینمی: \NSen{$<3\,\mathrm{g}/\mathrm{dL}$}.
\item هایپرکلسترولمی: \NSen{$>250\,\mathrm{mg}/\mathrm{dL}$}.
\item ادم.
\end{itemize}

\clearpage
\section*{سندرم نفروتیک ـ صفحهٔ ۲}
\subsection*{\NSen{Minimal Change Nephrotic Syndrome (MCNS)}}
\begin{itemize}
\item شایع‌ترین فرم \NSen{NS} اولیه در کودکان.
\item \NShl{مردان ۲ برابر زنان}.
\item پیک سنی: \NShl{قبل از مدرسه}.
\end{itemize}

\subsection*{\NSen{Focal Segmental Glomerulosclerosis (FSGS)}}
\begin{itemize}
\item شدت پروتئین‌اوری از \NSen{MCNS} کمتر است.
\item \NShl{بیش از ۲۵\% بیماران} به سمت \NShl{\NSen{ESRD}} می‌روند.
\end{itemize}

\subsection*{ممبرانوس نفروپاتی}
\begin{itemize}
\item نوجوانان، کودکان؛ عفونت مزمن (هپاتیت \NSen{B}، سفلیس، مالاریا).
\item مواجهه با برخی داروها (پنی‌سیلامین، طلا)، \NSen{SLE}.
\end{itemize}

\subsection*{\NSen{NS} مادرزادی و شیرخواران}
\begin{itemize}
\item \NSen{NS} مادرزادی: ظرف \NShl{۳ ماه نخست} تولد به وجود می‌آید؛ وراثت \NSen{AR}.
\item \NSen{NS} شیرخواران: ماه \NShl{۳ تا ۱۲}؛ ژنتیکی.
\end{itemize}

\medskip
\noindent
\NShl{در کودک با \NSen{NS}، سطح سرمی \NSen{$C_3\downarrow$}}
\(\longleftarrow\)
\NShl{علت \NSen{MCNS} نیست}.

\endgroup
\ifNephroticStandalone
  \expandafter\enddocument
\fi
